\documentclass[10pt]{article}
\usepackage[margin=0.78in]{geometry}
\usepackage{amsmath,amssymb,amsthm,booktabs,graphicx,microtype,placeins}
\usepackage[font=small,skip=3pt]{caption}
\usepackage[hidelinks]{hyperref}
\setlength{\parindent}{0pt}
\setlength{\parskip}{3pt plus 1pt minus 1pt}
\setlength{\textfloatsep}{7pt plus 2pt minus 2pt}
\setlength{\intextsep}{7pt plus 2pt minus 2pt}
\newtheorem{proposition}{Proposition}
\theoremstyle{definition}
\newtheorem{definition}[proposition]{Definition}
\title{\vspace{-1.7em}\textbf{A Semantic Construction Profile for}\\[-1pt]
\textbf{Exact Boolean Formula Complexity}}
\author{J. R. Landers}
\date{\vspace{-1.4em}}

\begin{document}
\maketitle
\vspace{-1.6em}

\begin{abstract}
Prior exhaustive synthesis showed that twelve elementary truth-table invariants
predict $77.7\%$ of exact four-input NAND/NOR formula-size variance and $84.3\%$
for AND/OR/NOT under strict descriptor holdout.  We investigate the residual
with 54 additional, input-renaming-invariant coordinates drawn from formula
and query complexity: decision trees, certificates, sensitivity, algebraic
normal form (ANF), signed Fourier structure, prime cubes, restrictions, and
hypercube topology.  The resulting 66-dimensional \emph{semantic construction
profile} forms 3,958 classes among all 65,536 functions, nearly matching the
3,984 input-permutation orbits.  A nonlinear regressor evaluated by five-fold
holdout of complete profile classes explains $95.7\%$, $94.7\%$, and $97.3\%$
of exact minimum-size variance in NAND, NOR, and AND/OR/NOT.  Conditional
permutation importance is dominated by decision-tree leaf count, followed by
certificate size, the Khrapchenko product, and linear ANF support.  Thus the
missing signal is structured: exact formula cost is most predictable from a
function's recursive decomposability, with boundary and algebraic structure
providing secondary information.  The result identifies ingredients for a
theory, but does not turn the learned relation into an exact formula or a new
asymptotic bound.
\end{abstract}

\section{From semantic correlation to construction potential}

\textbf{Takeaway.} The earlier descriptor measured coarse shape.  The expanded
profile measures ways a function can be decomposed, certified, or bounded.

For a gate basis $\mathcal L$, recall the exact tree-formula complexity
\[
 K_{\mathcal L}(f)=\min\{|e|:\operatorname{eval}_{\mathcal L}(e)=f\}.
 \tag{1}
\]
The target is known exactly for every $f:\{0,1\}^4\to\{0,1\}$ in NAND, NOR,
and ordered AND/OR/NOT by minimum-layer enumeration.  The original semantic
map $\psi(f)\in\mathbb R^{12}$ contained signed bias, four sorted influences,
five Walsh--Fourier energies by degree, algebraic degree, and input-permutation
stabilizer size.  It was compact, but it discarded how spectral
mass is signed, how ANF support is distributed, and how the truth table
decomposes under queries or restrictions.

\begin{definition}[Semantic construction profile]
Let
\[
 \Phi(f)=\bigl(\psi(f),D(f),C(f),S(f),A(f),W(f),P(f),R(f),H(f)\bigr)
 \in\mathbb R^{66},                                                   \tag{2}
\]
where $D,C,S,A,W,P,R,H$ respectively summarize exact decision trees,
certificates, local and higher-order sensitivity, ANF support, signed and
shape-sensitive Fourier data, prime cubes, restrictions and factorability,
and monochromatic hypercube topology.
\end{definition}

All coordinates are functions of the truth table; no minimum formula or gate
count is used.  Sorting or aggregation makes every coordinate invariant under
renaming inputs.  The term ``construction profile'' is operational: each
family records either a known obstruction to small formulas or a constructive
route to one.  Decision-tree leaves expand directly into a formula; certificate
complexity is linked to decision-tree complexity; the Khrapchenko product is a
De Morgan formula lower bound; prime implicants induce DNF/CNF upper bounds;
and random restrictions are a classical formula-size method
\cite{jukna2012,buhrman2002,impagliazzo1993,chandra1978}.

\begin{table}[!htbp]
\centering\small
\begin{tabular}{lrl}
\toprule
Profile family & Coordinates & Mathematical connection \\
\midrule
Original descriptor & 12 & influence, Fourier concentration, symmetry \\
Decision trees/certificates & 6 & constructive trees; query inequalities \\
Local/higher sensitivity & 9 & Khrapchenko and boundary complexity \\
ANF support & 5 & $\mathbb F_2$ gate composition \\
Fourier phase/shape & 17 & information lost by degree-wise energy \\
Prime cubes & 6 & canonical DNF/CNF constructions \\
Restrictions/factorability & 8 & shrinkage and recursive decompositions \\
Cube topology & 3 & fragmentation beyond boundary size \\
\bottomrule
\end{tabular}
\caption{The 66-coordinate profile.  The rows partition the coordinates; the
first row is the complete original descriptor.}
\label{tab:profile}
\end{table}

\section{Evaluation without semantic twins}

The expanded map partitions the complete universe into 3,958 exact profile
classes, compared with 675 original descriptor classes and 3,984 orbits under
input permutation.  Near-orbit resolution creates a leakage risk: a random
function split could place an identical profile in both train and test.
We therefore split \emph{profile classes}, not functions, into five folds.
Every class is represented once in fitting, weighted by its number of member
functions.  At evaluation, the within-class squared error is restored, so a
model receives no credit for variation that $\Phi$ cannot distinguish.

For fold $j$, a histogram gradient-boosting regressor $m_{-j}$ is trained on
the other four folds and evaluated by function-weighted
\[
 R^2_j=1-\frac{\sum_f(K_{\mathcal L}(f)-m_{-j}(\Phi(f)))^2}
 {\sum_f(K_{\mathcal L}(f)-\overline K_{\mathcal L,j})^2}.             \tag{3}
\]
This tests generalization to unseen semantic profiles.  It differs from the
earlier tie-aware $k$-nearest-neighbor experiment in both descriptor and model,
so the comparison is evidence of recovered signal rather than a one-variable
ablation.

Importance is computed only on held-out classes.  For coordinate $q$, we
permute $q$ among test profiles and report
\[
 I_q=\mathbb E_{j,r}\bigl[R^2_j-R^{2,\operatorname{perm}(q)}_{j,r}\bigr],
 \tag{4}
\]
over five folds and eight permutations per fold.  We also permute all
coordinates in a conceptual family jointly.  These are \emph{conditional}
importances: a low value may mean that correlated coordinates substitute for
one another, not that the underlying concept is irrelevant.

\section{Results}

\textbf{Takeaway.} The enriched profile closes most of the finite prediction
gap, and the learned dependence has a clear hierarchy.

\begin{table}[!htbp]
\centering\small
\begin{tabular}{lrrr}
\toprule
Language & Original strict $R^2$ & Expanded held-out $R^2$ & Fold s.d. \\
\midrule
NAND       & .777 & .957 & .005 \\
NOR        & .777 & .947 & .008 \\
AND/OR/NOT & .843 & .973 & .004 \\
\bottomrule
\end{tabular}
\caption{Prediction of exact minimum gate count over all 65,536 functions.
Both evaluations exclude exact descriptor twins; the expanded result uses
five-fold holdout of complete 66-feature profile classes.}
\label{tab:prediction}
\end{table}

The original descriptor ceiling was .893 for NAND/NOR and .931 for
AND/OR/NOT.  The expanded profile ceiling is .9996, .9996, and .9999: only
15, 14, and 2 profile classes, respectively, retain any target ambiguity.
The nonlinear model remains below those ceilings, showing that a small amount
of error is statistical/modeling error rather than missing truth-table signal.
Nevertheless, mean held-out $R^2$ rises by .180 for NAND, .169 for NOR, and
.130 for AND/OR/NOT relative to the original strict result.

\begin{figure}[!ht]
\centering
\includegraphics[width=0.98\textwidth]{importance_compact.pdf}
\caption{Conditional permutation importance on unseen profile classes.  Bars
average the three gate languages; points give language-specific values.
``New'' coordinates were absent from the original semantic descriptor.}
\label{fig:importance}
\end{figure}
\FloatBarrier

Decision-tree leaf count is the dominant individual coordinate, with mean
importance .283.  Mean certificate size follows at .097, then the Khrapchenko
product (.055) and the number of linear ANF monomials (.047).  Jointly,
decision-tree coordinates score .283, certificates .167, local sensitivity
.081, ANF support .068, prime cubes .051, and Fourier phase .039.  The original
Fourier-magnitude family remains useful (.023), while individual influences
become conditionally small (.002 jointly) because certificates, sensitivity,
and Khrapchenko coordinates encode overlapping boundary information.

Standalone association tells the same broad story without conditioning on the
regressor.  Decision-tree leaf count, the Khrapchenko product, and mean
certificate size each explain about $74\%$ of target variance under an adjusted
nonparametric correlation ratio.  No single coordinate suffices, however:
their conditional roles separate once competing semantic summaries are held
fixed.

\section{What the model appears to measure}

The hierarchy suggests the following finite structural hypothesis:
\[
 \text{formula cost}\ \approx\
 \text{recursive decomposition cost}
 +\text{boundary obstruction}
 +\text{algebraic/phase correction}.                                  \tag{5}
\]
The first term is witnessed by decision-tree leaves and certificates.  These
features ask how quickly restrictions force the output and therefore how many
semantic cases a tree formula must separate.  The second is witnessed by the
Khrapchenko product and local sensitivity, which measure the edge boundary
between zero and one inputs and have direct lower-bound precedent.  The third
distinguishes functions with similar boundary magnitude but different ANF
support or Fourier signs.

The correction is not merely a passive correlation.  XORing linear parity
into a function leaves all higher-degree ANF coefficients fixed.  Among 17,024
such pairs that also preserve the original descriptor exactly, NAND cost
changes by 1.271 gates on average and by as many as eight; matched non-parity
pairs average 1.167.  The analogous NOR values are 1.268 and 1.167.  For
AND/OR/NOT, however, parity-related pairs are \emph{closer} than controls
(.664 versus .813).  Phase sensitivity is therefore real but basis-relative;
``parity burden'' is not a universal scalar explanation.

Equation (5) is a research program rather than a closed form.  A theorem would
need explicit upper and lower inequalities connecting these quantities to
$K_{\mathcal L}$, while controlling their redundancy and basis dependence.
The present result supplies empirical targets for such inequalities: begin
with recursive restriction trees, add boundary lower bounds, then characterize
when algebraic phase makes those bounds loose.

\section{Toward an analytic statement}

The importance ordering points to a more precise theorem target.  Let $U(f)$
be the minimum number of leaves in a deterministic decision tree for $f$, and
let $B(f)$ be its Khrapchenko product (the product of the average numbers of
bichromatic hypercube edges incident to zero and one inputs).  A decision tree
can be compiled recursively by Shannon expansion into a formula, giving a
basis-dependent upper bound linear in $U(f)$.  Conversely, $B(f)$ lower-bounds
De Morgan formula leaf size and therefore, up to constant-factor translations,
the formula measures considered here \cite{jukna2012}.  The established
machinery thus has the schematic form
\[
 c_{\mathcal L}B(f)\ \le\ K_{\mathcal L}(f)+1\ \le\
 C_{\mathcal L}U(f),                                                   \tag{6}
\]
for suitable conventions and basis-dependent constants.

Our leading features sit at the two ends of (6): $U$ is first and $B$ third in
individual conditional importance.  Mean certificate size, ranked second,
describes how rapidly the upper construction terminates on typical inputs.
This suggests that the regressor is learning where $K_{\mathcal L}(f)$ lies
inside a classical lower--upper interval rather than discovering an unrelated
numerical coincidence.  ANF support and Fourier phase then predict when
boundary-based lower information or decision-tree upper information is loose.

Because the four-input universe is complete, this claim can be sharpened
without fitting another black-box predictor.  One can first compute the best
valid constants in (6) for each basis, then regress only the normalized
position inside the resulting interval.  Symbolic search can propose a small
certificate/phase correction, and exhaustive verification can either prove
the candidate for all four-input functions or return exact counterexamples.
The useful output would be an inequality or approximation guarantee with a
short semantic vector, not a 66-coordinate identity test.  Five-input tests
would then distinguish a genuine scaling law from a finite-universe fit.

\section{Limitations and next tests}

The universe is exhaustive but has only four inputs.  The profile is almost an
orbit identifier, so its near-unit ceiling does not itself constitute a compact
theory.  Gradient boosting and permutation importance describe predictive
dependence, not causation; importances can move among correlated features.
Finally, exact decision-tree quantities become expensive as $n$ grows.

The decisive next experiment is therefore not to add more four-input
coordinates.  It is to freeze a compact profile---decision-tree leaves,
certificate means, a Khrapchenko term, low-degree ANF support, and signed
Fourier counts---and test selected five-input families under the same
profile-class holdout.  In parallel, inequality search can test whether a
small combination sandwiches exact complexity within a basis-dependent
additive or multiplicative error.  Success would turn the observed semantic
construction potential into a mathematical theory; failure would localize the
remaining obstruction.

\paragraph{Reproducibility.}
All targets, 66 features, profile classes, folds, and permutation trials are
checked in with the companion experiment.  Five folds and eight permutations
per fold use fixed seeds; no cloud service or model API is required.

\begingroup
\footnotesize
\setlength{\parskip}{0pt}
\begin{thebibliography}{9}
\setlength{\itemsep}{0pt}
\bibitem{jukna2012}
S. Jukna, \emph{Boolean Function Complexity: Advances and Frontiers}.
Springer, 2012.
\bibitem{buhrman2002}
H. Buhrman and R. de Wolf, ``Complexity measures and decision tree
complexity: a survey,'' \emph{Theor. Comput. Sci.}, 288(1):21--43, 2002.
\bibitem{impagliazzo1993}
R. Impagliazzo, ``The effect of random restrictions on formula size,''
\emph{Random Struct. Algorithms}, 4(2):121--134, 1993.
\bibitem{chandra1978}
A. Chandra and G. Markowsky, ``On the number of prime implicants,''
\emph{Discrete Mathematics}, 24(1):7--11, 1978.
\bibitem{odonnell2014}
R. O'Donnell, \emph{Analysis of Boolean Functions}. Cambridge University
Press, 2014; revised arXiv edition, 2021.
\bibitem{friedman2001}
J. H. Friedman, ``Greedy function approximation: a gradient boosting
machine,'' \emph{Annals of Statistics}, 29(5):1189--1232, 2001.
\end{thebibliography}
\endgroup

\end{document}
